% Standalone response-letter draft for JMLR-25-3153-1.
% Compile twice with pdflatex. No external .bib, figures, or manuscript .aux
% files are required. Reviewer comments are blue; responses are black.
%
% IMPORTANT: The current non-experimental manuscript was not available when
% this letter was drafted. Red author actions, blank revision blocks, and
% unresolved manuscript numbers must be completed before submission.
% Non-experimental explanations below are proposed responses, not assertions
% that the corresponding manuscript edits or additional diagnostics exist.
%
% Evidence: experiments_body.tex, experiments_appendix.tex, data/*.csv,
% ../real_exps/*.csv, ../utility/exps.py, ../utility/res_rescaled.py,
% ../utility/cqhr.py, and the original submission and two referee reports.

\documentclass[11pt]{article}
\usepackage[margin=1in]{geometry}
\usepackage[T1]{fontenc}
\usepackage{lmodern}
\usepackage{amsmath,amssymb,bm}
\usepackage{booktabs,array}
\usepackage{enumitem}
\usepackage{xcolor}
\usepackage{microtype}
\usepackage[colorlinks=true,linkcolor=black,urlcolor=blue,citecolor=black]{hyperref}

\definecolor{ReviewerBlue}{RGB}{0,112,157}
\definecolor{AuthorRed}{RGB}{170,35,35}
\setlength{\parindent}{0pt}
\setlength{\parskip}{0.55em}
\setlength{\emergencystretch}{2em}
\setlist{topsep=0.4em,itemsep=0.3em,parsep=0pt}
\renewcommand{\thesection}{R\arabic{section}}
\renewcommand{\thesubsection}{\thesection.\arabic{subsection}}
\numberwithin{equation}{section}

\newcommand{\TSCP}{\textnormal{TSCP}}
\newcommand{\GWC}{\textnormal{TSCP-GWC}}
\newcommand{\LWC}{\textnormal{TSCP-LWC}}
\newcommand{\R}{\mathbb{R}}
\newcommand{\E}{\mathbb{E}}
\newcommand{\Prob}{\mathbb{P}}
\newcommand{\ind}{\mathbf{1}}
\newcommand{\pending}[1]{\textcolor{AuthorRed}{[INSERT: #1]}}
\newcommand{\authoraction}[1]{%
  \par\begingroup\small\color{AuthorRed}%
  \textbf{Author action (remove before submission).} #1\par\endgroup}
\newcommand{\reviewercomment}[2]{%
  \par\medskip\noindent\textbf{#1}\par\nopagebreak
  \begingroup\color{ReviewerBlue}\itshape #2\par\endgroup
  \noindent\textbf{Response.} }
\newcommand{\revisionblock}[2]{%
  \par\begingroup\color{AuthorRed}\small
  \textbf{Manuscript revision to insert (#1).} #2\par
  \textbf{Exact revised text or description of the change:}\par
  \vspace{0.35em}\hrule height 0.3pt\vspace{1.4em}
  \hrule height 0.3pt\vspace{0.6em}
  \textbf{Location in the revised manuscript:}
  Section/Appendix \rule{1.3cm}{0.3pt},
  page \rule{1cm}{0.3pt},
  lines \rule{1.3cm}{0.3pt};
  theorem/equation/figure \rule{1.4cm}{0.3pt}.
  \par\endgroup}

% Fill each manuscript reference below once, after assembling the paper.
% Comments give the verified labels in the experiment TeX files.
% Do not infer revised numbering from the old manuscript.
\newcommand{\SecExperiments}{Section~\pending{EXP}} % sec:numerical-experiments
\newcommand{\SecMetrics}{Section~\pending{METRICS}} % sec:experiment-methods
\newcommand{\SecAbsolute}{Section~\pending{ABS}} % sec:absolute-residual-experiments
\newcommand{\SecCQR}{Section~\pending{CQR}} % sec:cqr-experiments
\newcommand{\SecReal}{Section~\pending{REAL}} % sec:real-data-experiments
\newcommand{\AppAdditional}{Appendix~\pending{ADDITIONAL}} % app:additional-experiments
\newcommand{\AppUncertainty}{Appendix~\pending{UNCERTAINTY}} % app:coverage-uncertainty
\newcommand{\AppDependence}{Appendix~\pending{DEPENDENCE}} % app:dependence-sensitivity
\newcommand{\AppHeterogeneity}{Appendix~\pending{HETEROGENEITY}} % app:heterogeneity-sensitivity
\newcommand{\AppSmall}{Appendix~\pending{SMALL}} % app:alpha-small-sample
\newcommand{\AppHeavy}{Appendix~\pending{HEAVY}} % app:heavy-tails
\newcommand{\AppContamination}{Appendix~\pending{CONTAMINATION}} % app:contamination
\newcommand{\AppCQR}{Appendix~\pending{CQR-SENSITIVITY}} % app:cqr-sensitivity
\newcommand{\AppShape}{Appendix~\pending{SHAPE}} % app:shape-template
\newcommand{\FigAbsolute}{Figure~\pending{ABS-OVERVIEW}} % fig:body-abs-overview
\newcommand{\FigAbsoluteBars}{Figure~\pending{ABS-BARS}} % fig:body-abs-coordinate-bars
\newcommand{\FigHeteroskedastic}{Figure~\pending{PARTIAL-HETERO}} % fig:body-partial-hetero
\newcommand{\FigCQR}{Figure~\pending{CQR-COMPARISON}} % fig:body-cqr
\newcommand{\FigCQBars}{Figure~\pending{CQR-BARS}} % fig:body-cqr-coordinate-bars
\newcommand{\FigRuntime}{Figure~\pending{REAL-RUNTIME}} % fig:body-real-runtime
\newcommand{\TabReal}{Table~\pending{REAL-DATA}} % tab:real-data
\newcommand{\FigHeavy}{Figure~\pending{HEAVY-TAILS}} % fig:app-heavy-tails
\newcommand{\FigContamination}{Figure~\pending{CONTAMINATION}} % fig:app-contamination
\newcommand{\FigSmall}{Figure~\pending{SMALL-CALIBRATION}} % fig:app-small-calibration
\newcommand{\FigBase}{Figure~\pending{CQR-BASE}} % fig:app-cqr-base
\newcommand{\FigShift}{Figure~\pending{CQR-SHIFT}} % fig:app-cqr-shift
\newcommand{\FigShape}{Figure~\pending{SHAPE-TEMPLATE}} % fig:app-shape-template

\begin{document}

\begin{center}
  {\Large\bfseries Revision of\par}
  \vspace{0.4em}
  {\Large\bfseries ``Interpretable Multivariate Conformal Prediction\par
  with Fast Transductive Standardization''\par}
  \vspace{0.8em}
  Manuscript ID: JMLR-25-3153-1\par
  Yunjie Fan and Matteo Sesia\par
  \pending{resubmission date}\par
  \vspace{0.8em}
  {\large\bfseries Response to the Editor and Reviewers}
\end{center}

We thank the Editor and both reviewers for their careful reading and
constructive feedback. The comments identify important opportunities to
clarify the scope of the theoretical results, strengthen the empirical
comparisons, and make the construction easier to follow. We have expanded
and reorganized the numerical experiments to address these points. Below
we summarize the experimental changes and provide a point-by-point response
to the decision letter and each report. Reviewer comments are reproduced
in blue italics, followed by our responses in black. Section and figure
references refer to the revised manuscript unless explicitly identified as
references to the original submission.

\authoraction{This is a working response letter, not a submission-ready
claim that every concern has already been resolved. Complete all red fields
and revision blocks using the current manuscript. In particular:
(i) reconcile the moment and tie assumptions with the actual revised
theorems; (ii) obtain the real-data regularity/fallback and backward-search
diagnostics; (iii) supply the requested real-data coordinate-wise evidence
or explicitly explain its omission; (iv) verify revisions to the literature
review, schematic figures, and vector notation; and (v) confirm that all
modified manuscript sections are highlighted. The experiment files and
saved results support the numerical statements below, but they do not
contain those missing diagnostics or the current non-experimental text.
Replace the centralized reference macros in the preamble after final
manuscript numbering is known.}

\section{Summary of Main Changes}
\label{sec:summary}

\paragraph{Dimension-specific comparisons in the main text.}
We have added coordinate-wise bar plots for a ten-dimensional
absolute-residual experiment and a three-dimensional quantile-residual
experiment (\FigAbsoluteBars{} and \FigCQBars{}). These report full
outcome-space interval lengths and are interpreted together with joint
coverage, making the effect of coordinate adaptation visible beyond an
aggregate volume comparison.

\paragraph{Quantile-regression experiments and additional baselines.}
The main text now compares TSCP with native CQHR using common fitted
quantile regressions (\SecCQR{}). The TSCP comparison uses capped
nonnegative CQR scores, while methods that accept signed scores retain
their native signed-score constructions. We also report sensitivity to
the base-interval miscoverage and to a constant-shift transformation in
\AppCQR{}. A separate low-dimensional comparison with a hyperrectangular
shape-template implementation based on Tumu et al.\ \cite{tumu2024} is
included in \AppShape{}.

\paragraph{Dependence and partial heteroskedasticity.}
The main absolute-residual experiment now includes correlated,
heterogeneous Gaussian noise. An additional main-text experiment makes
only five of ten coordinates heteroskedastic in an input feature
(\FigHeteroskedastic{}). Further experiments vary the strength of
dependence, the severity of coordinate heterogeneity, the target coverage,
and the calibration sample size (\AppDependence{}, \AppHeterogeneity{},
and \AppSmall{}).

\paragraph{Wall-clock evidence in the main text.}
We have added conformal-region construction times to the main synthetic
comparison and a dedicated runtime figure for the six retained real-data
benchmarks (\FigAbsolute{} and \FigRuntime{}). These measurements exclude
the common regression fit and residual computation. They supplement,
rather than replace, the worst-case complexity analysis.

\paragraph{Heavy tails, contamination, and limitations.}
\AppHeavy{} explicitly identifies the infinite-variance cases as stress
tests outside the finite-second-moment assumptions of the original
submission. \AppContamination{} adds an exchangeable contamination study
that distinguishes near-nominal coverage from severe inflation of region
volume. These results do not establish robust efficiency or an extension
of a theorem beyond its stated assumptions.

\paragraph{Reproducibility and uncertainty.}
Each revised synthetic repetition redraws training, calibration, and test
observations independently, with the underlying regression coefficients
held fixed. Absolute-residual studies use 200 repetitions, 7,200 training
observations, and 800 test observations. CQR studies use 30 repetitions,
2,400 training observations, and 600 test observations, consistently across
the main comparison and its sensitivity analyses. The low-dimensional
shape-template study uses 100 repetitions and the same 2,400/600
training/test sizes. Real-data results retain 200 random splits of each
fixed dataset. Across-repetition uncertainty is retained in the result
files and discussed in \AppUncertainty{}.

\revisionblock{summary of non-experimental revisions}{Insert a short
summary of the actual revisions to the introduction, assumptions,
methodological explanation, schematic figures, and notation. Do not
state that a theoretical extension or additional diagnostic has been
added unless it appears in the revised submission.}

\section{Response to the Editor}
\label{sec:editor}

\reviewercomment{E1. Completeness of the revision and response}{%
The manuscript was sent to two expert Reviewers, who each noted several
positive aspects of the work while also highlighting areas in which it
could be improved. One Reviewer felt that the empirical assessment was
incomplete and suggested additional results and baselines which could
strengthen the work, while the other Reviewer was more positive about the
methodology and empirical assessment but expressed some difficulty in
digesting elements of the manuscript, suggesting areas for presentational
improvement. I would therefore like to invite a resubmission which
addresses the two Reviewers' feedback in full, together with a
point-by-point response to the Reviewers' reports.}

We appreciate the opportunity to revise and resubmit. The new empirical
material addresses the requests for coordinate-specific comparisons,
CQHR, a shape-template baseline, partial heteroskedasticity, outlier
contamination, and main-text runtime evidence. We have also separated
empirical observations from the scope of the formal guarantees. The
responses below address the methodological and presentation questions
individually, including the limitations of the new comparisons rather
than interpreting them as universal dominance results.

\revisionblock{E1: final revision overview}{Insert the final description
of the non-experimental changes and confirm that the diagnostic and
presentation requests flagged in this letter have been completed or
explicitly addressed.}

\reviewercomment{E2. Highlighting modified manuscript sections}{%
For the convenience of the Reviewers, I kindly request that modified
sections of the manuscript are indicated using a different colour font.}

Thank you for this instruction. The replacement experimental sections in
the main text and supplement are marked in blue. The corresponding
highlighting of changes elsewhere in the manuscript is
\pending{confirm completion and identify the marked manuscript version}.

\authoraction{Both experiment TeX files already enclose their replacement
text in a blue color group. The current introduction, methodology, proofs,
and other non-experimental sections were not available for inspection.
Do not claim that all changes are highlighted until the assembled
manuscript has been checked.}

\section{Detailed Response to Reviewer 1}
\label{sec:reviewer1}

We thank the reviewer for the positive assessment of the problem,
technical contribution, and progression from the oracle construction to
the practical method. We particularly appreciate the emphasis on clearer
exposition and a more careful distinction between validity, computational
cost, and robustness.

\subsection{Main Concerns}

\reviewercomment{R1.C1. Intuition for the local construction and rectangular enclosure}{%
The paper is fairly dense, and the methodological core is harder to parse
than it needs to be. The high-level intuition is good, but once the paper
gets into the transductive oracle, link functions, worst-case
constructions, and the indexing used in the local partition argument,
the presentation becomes difficult to follow. I think the paper would
benefit from more intuition in Section 3, especially around why the local
worst-case construction is the ``right'' relaxation and how the
rectangular enclosure affects efficiency in practice.}

We agree that the role of each step should be clear before introducing
its detailed notation. The intended explanation is as follows.

\paragraph{Why use a transductive oracle?}
Estimating coordinate means and scales from the calibration residuals
alone and then treating the resulting transformation as fixed does not
generally preserve the symmetry between calibration and test scores.
The oracle instead includes the test residual in the standardization.
Its statistics are symmetric in all $n+1$ residual vectors, which explains
the conformal rank argument. The obstacle is computational, since the
test residual is unknown, not a need to estimate a parametric dependence
model.

\paragraph{What do the link functions accomplish?}
The link functions translate a scalar standardized-score threshold into
coordinate-wise residual bounds using observable calibration statistics.
They separate two tasks: inverting the transductive score inequality and
finding a computable conservative threshold. This is why the subsequent
discussion can focus on bounding an unknown oracle quantile rather than
searching directly over every possible test outcome.

\paragraph{Why localize a worst-case bound?}
The global construction protects against every admissible test residual
at once. This can be overly conservative because the values producing
the most unfavorable bounds need not be relevant when the test residual
is restricted to a particular region. A local construction uses the
working hypothesis that the test residual lies in one cell of a
data-driven partition. It then constructs an upper bound valid under
that hypothesis. The union over cells retains the relevant oracle
acceptance event, while its restriction to the global region controls
the outer extent. Thus, ``local'' refers to a localization of the
worst-case calculation in residual/outcome space; it is not a claim of
conditional coverage at a fixed covariate value.

\paragraph{What is lost by enclosing the union in a rectangle?}
The rectangular enclosure may include points that lie in none of the
local sets. This can only enlarge the reported region, so it trades
some efficiency for an interpretable Cartesian product. The construction
is a tractable conservative relaxation, not a claim to solve an
unrestricted minimum-volume optimization problem. In particular,
coordinate-wise endpoint calculations can be conservative if they use
local feasibility information separately. We should not describe the
reported rectangle as the exact smallest enclosure without establishing
that stronger property.

The new coordinate-length plots in \FigAbsoluteBars{} and \FigCQBars{}
show the informativeness of the final reported rectangle, while the
joint coverage and volume comparisons show its overall empirical
behavior. These plots do not directly estimate the volume lost solely
in the enclosure step. Any quantitative claim about that specific
approximation requires a separate comparison with the unenclosed local
union, rather than a comparison only with other conformal methods.

\revisionblock{R1.C1}{Insert the actual new roadmap and intuitive
explanation in Section 3, and the revised explanation of the link
functions and enclosure. If the revised paper retains an oracle/local
union/enclosure comparison from the old supplement, insert its final
location and the exact supported conclusion. Otherwise retain the
qualitative explanation without claiming a new enclosure-gap study.}

\reviewercomment{R1.C2. Assumptions, infinite-variance experiments, and negative residuals}{%
More discussion of the assumptions and how they relate to the experiments
are needed. The theory assumes nonnegative residuals with finite first
and second moments, but the appendix also emphasizes robustness under
heavy-tailed settings where the population variance is infinite. That
may be empirically true, but it creates some tension with the formal
assumptions, and the paper should be more explicit about what exactly is
proved versus what is only observed experimentally. Along similar
lines, the treatment of negative residuals is deferred to workarounds
and future work, which is understandable, but it does narrow the
immediate scope of the method.}

We agree that the previous presentation did not distinguish these issues
sharply enough. There are three different questions: whether a conformal
rank argument is valid, whether the particular TSCP formulas and
containment arguments are justified under the stated assumptions, and
whether the resulting region is statistically efficient.

\paragraph{Finite moments and the scope of the theorem.}
The standard exchangeability-based rank argument does not itself require
population means or variances. However, that observation alone does not
remove assumptions from every subsequent result in the paper. In the
original submission, Assumption 1 included finite positive variance and
was explicitly invoked by the final TSCP validity theorem, as well as
by intermediate results. Therefore, the original theorem should not be
cited as already proving validity for infinite-variance residuals.
The fact that an observed finite sample has finite numerical entries
is also not a substitute for a population moment assumption.

The revised discussion must identify the assumptions of each actual
theorem. If an assumption has been weakened, the corresponding revised
statement and proof must be supplied; otherwise the heavy-tail results
remain empirical evidence outside that theorem's scope.

\paragraph{What the heavy-tail experiment establishes.}
\AppHeavy{} and \FigHeavy{} consider independent, common-scale Student
$t_\nu$ noise with $d=10$ and $\nu\in\{1.5,2,3\}$. The first two cases
have infinite variance; the third has finite variance. TSCP's observed
coverage remains near $0.9$, but its volume can be very large. For
example, at $n=30$ the mean normalized volume is approximately
$1.59\times10^{20}$ for $\nu=1.5$, compared with $5.92\times10^{13}$
for $\nu=3$. At $n=500$, the corresponding values are approximately
$8.42\times10^{13}$ and $6.07\times10^7$. These results are a stress
test, not evidence that the procedure is efficiently robust to arbitrarily
heavy tails or that a finite-moment theorem applies unchanged.

\paragraph{What the quantile-residual experiment establishes.}
To make the treatment of signed scores concrete, \SecCQR{} now uses
the CQR score
\begin{equation}
  S_j(x,\bm y)
  =\max\{\widehat q_{\ell,j}(x)-y_j,
           y_j-\widehat q_{u,j}(x)\}
\end{equation}
and supplies TSCP with its positive part $E_j=\max\{S_j,0\}$. We also
study $S_j+C$ in \AppCQR{}. Capping satisfies nonnegativity but prevents
the final interval from shrinking the base interval. Moreover, it
introduces an atom at zero, so it does not satisfy the no-point-mass
part of the original Assumption 1 verbatim. The scope of the revised
theory must therefore address ties and possible degenerate empirical
scales explicitly, rather than treating nonnegativity as the only
condition. Our detailed response to R2.M7 below explains this distinction.

\revisionblock{R1.C2}{Insert the exact revised assumption(s), theorem
statement(s), and explanatory paragraph separating formal validity from
empirical heavy-tail performance. Include the treatment of capped-score
ties and zero empirical variance, or explicitly delimit this experiment
as an empirical application not covered verbatim by the original theorem.}

\authoraction{Check the current non-experimental text before choosing
the final wording. The existing experimental appendix describes the
infinite-variance cases as outside finite-second-moment ``efficiency''
assumptions, but the old submission also invokes Assumption 1 in the
validity theorem. Align that appendix wording with the actual revised
theorems; this letter deliberately does not assert an unseen extension.}

\reviewercomment{R1.C3. Quadratic worst-case complexity and wall-clock evidence}{%
While the complexity discussion is appreciated, I am not fully convinced
by the computational story yet. The worst-case complexity is still
quadratic in the calibration size up to dimension factors, and it would
help to see more direct runtime comparisons in the main paper, not only
prediction-set volume comparisons. Since one of the selling points is
tractability, clearer wall-clock evidence would strengthen the paper.}

Thank you for this suggestion. We have added construction times to the
main synthetic figure (\FigAbsolute{}) and a dedicated real-data
runtime comparison (\FigRuntime{}). On the six retained real-data
benchmarks, TSCP's mean conformal-region construction time ranges from
approximately $0.74$ to $7.46$ milliseconds. The tested calibration
sizes range from $15$ to $490$, and the output dimensions range from
$2$ to $16$. TSCP is slower than the simpler Unscaled Max and Point
CHR implementations. On the two 16-dimensional SCM datasets it is
substantially faster than the empirical-copula implementation, whose
mean times are approximately $95$--$118$ milliseconds.

These are wall-clock times for constructing a region from already
computed calibration residuals. They do not include fitting the common
predictive model or computing its residuals. The real-data figure uses
the retained timings from the original real-data runs; it is not a new
large-scale benchmark. The synthetic runtime figure uses the revised
fresh-sample experiments.

We agree that these measurements do not eliminate the worst-case
quadratic dependence. In the original analysis, constructing the
coordinate thresholds costs $O(d^2n\log n)$ when binary searches are
available throughout, and $O(d^2n^2)$ in the worst case. Once computed,
the thresholds are reused across test inputs; applying them to $m$
test points adds $O(dm)$ arithmetic for the usual residual choices,
excluding evaluation of the fitted predictor. The defensible empirical
conclusion is that construction is inexpensive at the tested sizes,
not that the algorithm has subquadratic worst-case complexity or has
been benchmarked at arbitrarily large $n$. Branch-specific evidence is
addressed separately in R1.Q2.

\revisionblock{R1.C3}{Insert the actual revised complexity discussion
and its main-text pointers to the runtime figures. Confirm the bounds
against the current algorithm. Add the hardware/software timing context
from the original runs where available; do not substitute the current
machine specification for an unverified historical environment.}

\reviewercomment{R1.C4. Contamination and robust standardization}{%
The paper's robustness to outliers seems mixed. The authors acknowledge
this in the discussion and point to rf2 as an example where influential
observations can substantially enlarge the sets. That honesty is
appreciated, but it also suggests that the current standardization
scheme may be brittle in some realistic settings. I would have liked
either a stronger empirical stress test on contamination/outliers or a
small robust variant, even if only heuristic.}

We agree. We have followed the suggestion to add a stronger stress test,
rather than presenting an unanalysed robust modification as if it had
the same guarantees as the original construction.

\AppContamination{} and \FigContamination{} use $d=10$, $n=100$, and
200 repetitions. Independently for each observation, all ten Gaussian
noise coordinates are multiplied by $10$ with probability
$\varepsilon\in\{0,0.01,0.05,0.10\}$. The same mixture mechanism is
used for training, calibration, and test observations. Thus the study
preserves the common data-generating distribution across these samples;
it is not an adversarial-corruption or train/test distribution-shift
experiment.

The results separate coverage from informativeness. TSCP's observed
coverage is approximately $0.904$ without contamination and $0.902$
at $10\%$ contamination. Its median normalized volume, however,
increases from approximately $5.38\times10^{10}$ to
$1.10\times10^{17}$. We report median volume in this stress test to
avoid having the graphical summary determined entirely by a few extreme
regions; the choice of summary is stated explicitly. Point CHR and
Unscaled Max also exhibit substantial inflation. These findings support
the reviewer's concern: near-nominal coverage does not imply robust
efficiency of mean/standard-deviation scaling.

Replacing those statistics by, for example, robust location and scale
estimators is an interesting direction, but it would require revisiting
the link functions, worst-case bounds, and computational shortcuts.
We do not claim to have implemented or proved such a variant. The
retained real-data table also continues to show that Point CHR has
smaller average volume than TSCP on \texttt{rf2}; the revision does not
hide this unfavorable case.

\revisionblock{R1.C4}{Insert the revised discussion of outlier
sensitivity, the reference to the new contamination experiment, and
the appropriately qualified future-work statement. Avoid describing
stable coverage under exchangeable contamination as resistance to
arbitrary contamination of the calibration sample alone.}

\subsection{Additional Questions}

\reviewercomment{R1.Q1. Formal assumptions versus infinite-variance experiments}{%
Can the authors clarify more explicitly the gap between the formal
assumptions used in the validity results and the heavy-tailed
experiments where variance may be infinite?}

Yes. As explained in R1.C2, an exchangeability-based conformal rank
argument and the assumptions of the full TSCP theorem are distinct.
The original theorem explicitly included Assumption 1; consequently,
the infinite-variance experiments do not establish its applicability
outside that assumption. We identify $t_{1.5}$ and $t_2$ as the
infinite-variance cases and $t_3$ as the finite-variance comparison in
\AppHeavy{}. Any stronger theoretical statement in the revision must
be justified by the revised theorem and proof, not by these simulations.

\authoraction{After completing the R1.C2 revision block, insert its
final assumption/theorem and discussion references here:
\pending{assumption/theorem numbers and manuscript pages}. This answer
should use precisely the same scope as the main response above.}

\reviewercomment{R1.Q2. Frequency and cost of backward search}{%
How often does the algorithm fall into the less favorable
backward-search regime in practice, and what are the actual runtime
implications on the datasets considered?}

This is an important diagnostic that is not identified by an overall
mean runtime alone. It is also different from the failure of the
regularity condition discussed in R2.M3. If the mean-index cell is
empty after restriction to the global region, the algorithm falls back
to GWC. When that cell is nonempty, some coordinates may still require
backward search instead of binary search.

The existing real-data result files contain mean construction times,
coverage, and volume, but do not record these branch indicators. Thus
they cannot support a numerical claim that backward search was rare
or absent. The main-text runtime figure answers the question of overall
observed cost, but not the requested frequency or conditional cost of
the backward-search regime.

For each dataset, the appropriate additional diagnostic distinguishes
the following quantities over its $B$ splits:
\begin{align}
  \widehat p_{\mathrm{fallback}}
  &=\frac{1}{B}\sum_{b=1}^{B}
    \ind\{\mathcal Y_{\bm h^*,b}=\varnothing\},\\
  \widehat p_{\mathrm{backward,run}}
  &=\frac{1}{B}\sum_{b=1}^{B}
    \ind\{\text{run $b$ uses backward search in at least one coordinate}\}.
\end{align}
It is also useful to report the fraction of coordinate searches that
are backward among runs that do not fall back, and construction times
conditional on the branch. The denominator must be specified in each
case; these are not interchangeable frequencies.

\revisionblock{R1.Q2: measured branch diagnostics}{Insert the measured
per-dataset counts and percentages for the six real-data benchmarks,
the denominator for each rate, and the mean or median construction
times for binary-only, backward-search, and fallback runs. State
whether these were collected from a replay of the original splits
or from new diagnostic runs. Include the table/figure location. If
the measurements remain unavailable, retain the limitation explicitly
instead of claiming this request has been fully resolved.}

\authoraction{This is a substantive unresolved empirical request, not
just a missing manuscript citation. Instrument the actual branch
conditions in \texttt{standardized\_prediction}; do not infer the
branch from a small runtime or from equality of the returned GWC and
local rectangles. Conditional runtime for a branch with zero observed
runs should be reported as not applicable, not as zero milliseconds.}

\reviewercomment{R1.Q3. Gains with quantile-regression residuals}{%
Do the authors expect the gains of TSCP to persist when using
quantile-regression residuals rather than absolute mean-regression
residuals?}

The new CQR experiments provide evidence that coordinate adaptation can
remain useful, although the magnitude of the improvement depends on
the fitted base intervals. We use three Gaussian outcomes with noise
scales $(3,2,1)$, the same independently trained quantile models for
all methods in a repetition, and final target coverage $0.9$. The
base-interval miscoverage is $0.5$. TSCP uses capped scores, Unscaled
Max and Empirical Copula use raw signed scores, and CQHR uses its
native construction. This distinction is now explicit in the text
and figure captions.

At $n=100$, TSCP has average coverage $0.9045$ and mean full
outcome-space volume approximately $2.91\times10^4$, compared with
$0.8973$ and $3.43\times10^4$ for Unscaled Max. The coordinate-wise
mean full lengths are $(36.2,36.9,20.9)$ for TSCP and
$(33.3,33.8,29.0)$ for Unscaled Max. Thus the observed gain is not
that every coordinate becomes shorter: TSCP uses more width in the
first two coordinates and substantially less in the third. The CQHR
comparison and its limitations are discussed in R2.M4.

Across the four displayed calibration sizes, TSCP's observed coverage
ranges from approximately $0.899$ to $0.922$. The small shortfall
at $n=50$ is accompanied by a repetition-level standard deviation
of approximately $0.0457$, rather than being presented as exact
finite-sample calibration. \AppCQR{} additionally varies the base
miscoverage and examines the shifted-score alternative. The evidence
supports usefulness in the tested setting, not a universal gain for
all quantile models, all score transformations, or all coordinates.

\section{Detailed Response to Reviewer 2}
\label{sec:reviewer2}

We thank the reviewer for the careful assessment of the construction
and for the specific suggestions on baselines, score transformations,
dimension-specific behavior, and presentation. These comments led to
substantial additions to the experiments and several important
clarifications of what is and is not implied by the method.

\subsection{Major Comments}

\reviewercomment{R2.M1. Point CHR versus CQHR and additional data splitting}{%
In section 1.2 you discuss that Sampson and Chan (2024) require an
additional independent dataset. This is true for their Point CHR
method, but not for their quantile regression approach (CQHR).}

The reviewer is correct, and the distinction should be made explicitly.
Point CHR uses a separate calibration stage to determine coordinate
scales before final conformal calibration. CQHR instead derives its
relative coordinate adjustments from fitted quantile-interval widths
and does not require that additional scale-estimation split. It still
uses the ordinary separation between model training and conformal
calibration. This is the distinction between the two procedures of
Sampson and Chan\ \cite{sampson2024}.

Accordingly, the data-efficiency motivation for TSCP should be stated
relative to procedures that estimate residual location/scale using
an auxiliary split, rather than attributed to every method in that
paper. The new main-text methods description makes the Point CHR/CQHR
distinction, and \SecCQR{} includes a native CQHR comparison.

\revisionblock{R2.M1}{Insert the corrected literature-review text in
the revised counterpart of Section 1.2, with its exact location. Check
that the abstract, introduction, and discussion do not repeat the
overgeneralized claim about every method of Sampson and Chan.}

\reviewercomment{R2.M2. Dimension-specific comparisons}{%
After Theorem 1 you discuss that the overall coverage cannot be used
to determine if the dimension-specific coverage is tight. However,
both the real data analyses and simulation studies lack
dimension-specific comparisons (e.g., comparing prediction interval
lengths or listing marginal coverages across the 10 dimensions in
the first simulation study for TSCP). Adding these to your numerical
comparisons would improve them greatly.}

We agree that overall coverage and a product of interval lengths do
not adequately describe which coordinates are informative.

\paragraph{Absolute-residual comparison.}
\FigAbsoluteBars{} reports coordinate-wise mean full lengths in the
ten-dimensional Gaussian experiment with scales $(10,9,\ldots,1)$,
equicorrelation $\rho=0.6$, and $n=100$. TSCP's mean length decreases
from approximately $49.4$ in the noisiest coordinate to $5.0$ in the
quietest, whereas Unscaled Max assigns approximately $38.0$ to every
coordinate. The corresponding joint coverages are approximately
$0.907$ and $0.902$. Empirical Copula and Point CHR are included in
the same plot; the shorter copula intervals must be interpreted
alongside its joint coverage of approximately $0.873$.

\paragraph{Quantile-residual comparison.}
\FigCQBars{} reports the analogous three-dimensional comparison at
$n=100$. TSCP's mean full lengths are $(36.2,36.9,20.9)$, compared
with $(33.3,33.8,29.0)$ for Unscaled Max and $(46.2,120.5,66.9)$
for CQHR. These are lengths in the original outcome space, averaged
over test observations and repetitions, not thresholds in transformed
score space. The main text also contains a coordinate profile for
the partial-heteroskedasticity experiment.

\paragraph{Interpretation and the real-data part of the request.}
These plots document coordinate-specific width allocation, not equal
marginal coverage across coordinates. A joint coverage lower bound
implies a lower bound for each coordinate's marginal coverage, but
does not imply that those marginal coverages are equal or close to
the joint target. The saved real-data tables contain only aggregate
coverage, volume, and runtime, so the new synthetic bar plots do not
by themselves answer the real-data component of this comment.

\revisionblock{R2.M2: real-data coordinate comparisons}{Insert the
additional real-data coordinate-wise length or marginal-coverage
results, identifying the dataset(s), methods, units, number of splits,
and table/figure location. If these results are not added, explain
that limitation explicitly rather than saying that both parts of
the request have been completed. Also insert any revised wording
after Theorem 1 that distinguishes joint validity, marginal coverage,
and coordinate-wise informativeness.}

\reviewercomment{R2.M3. Frequency of failure of the regularity condition}{%
In section 3.4.4 it is stated: ``Under the regularity condition
$\mathcal Y_{\bm h^*}\ne\varnothing$, which can be immediately
verified and typically holds in practice...'' How often did this
not hold in the real data experiments, if at all?}

Thank you for requesting a direct measurement. The condition concerns
the mean-index cell after its restriction to the global region. If
the condition fails, the implemented shortcut returns the GWC
rectangle rather than applying the local search. This fallback is
part of the construction; the practical consequence can be a more
conservative region, not an undefined result. Any validity statement
still uses the assumptions of the relevant theorem.

However, the saved real-data summaries do not contain a fallback
indicator, so they cannot establish that the condition always held.
We should not substitute the qualitative phrase ``typically holds''
for the requested counts. The relevant reporting unit is a
calibration split: the condition is checked once for the common
calibration residuals, not independently for each test observation.
The diagnostic should report, for every dataset, the number of
fallbacks out of the 200 splits and the resulting percentage.

\revisionblock{R2.M3: regularity-condition audit}{Insert the six
measured fallback counts and percentages, together with the revised
table location and the corresponding correction or qualification
of ``typically holds.'' Reuse the same diagnostic run described
in R1.Q2, but keep fallback frequency separate from backward-search
frequency. Do not fill an unmeasured count with zero.}

\authoraction{In the old submission, the sentence beginning ``The
special case'' in Section 3.4.4 appears to use a nonempty-set symbol
while describing the exceptional empty-cell case. Check the current
version for this possible sign/notation error when inserting the
actual revised passage.}

\reviewercomment{R2.M4. Missing CQHR baseline}{%
The numerical studies do not include a comparison with CQHR from
Sampson and Chan (2024), which performed better than Point CHR in
their paper.}

We have added CQHR as a main-text baseline in \SecCQR{}, rather than
using Point CHR as a substitute for it. All methods in this comparison
share the same fitted lower and upper quantile regressions and the
same training, calibration, and test observations within a repetition.
Each coordinate's quantiles are fitted with gradient boosting using
100 stages, depth $5$, learning rate $0.05$, and minimum leaf size $5$.
We use 2,400 training observations, 600 test observations, and 30
independent repetitions at $n\in\{30,50,100,200\}$.

The CQHR implementation follows the native interval-width rescaling
construction, using the first coordinate as reference. It operates
on signed CQR scores rather than on TSCP's capped transformation.
Predicted lower and upper limits are ordered pointwise before
computing scores for all methods in the comparison.
The code uses a $10^{-12}$ numerical floor when forming ratios of
base-interval lengths. This implementation detail is relevant when
interpreting very small fitted widths and is not an additional
calibration split.

At $n=100$, the principal results are:
\begin{center}
\begin{tabular}{lrr}
\toprule
Method & Mean joint coverage & Mean outcome-space volume\\
\midrule
Unscaled Max (signed scores) & $0.8973$ & $3.43\times10^4$\\
CQHR (native construction) & $0.8936$ & $1.29\times10^8$\\
TSCP (capped scores) & $0.9045$ & $2.91\times10^4$\\
\bottomrule
\end{tabular}
\end{center}
The corresponding repetition-level coverage standard deviations are
approximately $0.0340$, $0.0318$, and $0.0374$. The table is therefore
not intended to characterize a small empirical shortfall as evidence
against a conformal coverage guarantee.

The large CQHR mean volume is also highly variable: its
across-repetition volume standard deviation at this setting is
approximately $3.79\times10^8$. We report the observed result without
interpreting it as a general deficiency of CQHR. It is specific to
the fitted quantile models, width-ratio construction, and reference
coordinate in this design. The coordinate-wise comparison is in
\FigCQBars{}, and \FigBase{} varies the base-interval miscoverage
over $0.3,0.5,0.7,0.9$. Those checks help characterize the comparison,
but do not constitute a comprehensive tuning study over quantile
models or reference coordinates.

\authoraction{Confirm that the final experiment implementation notes
include the reference-coordinate choice and numerical length floor.
The current experiment narrative already states the model settings
and the restricted interpretation of CQHR's performance. Do not
claim that the cause of every extreme CQHR volume has been diagnosed
without an additional width-ratio audit.}

\reviewercomment{R2.M5. Missing convex shape-template baseline and related work}{%
The numerical studies and lit review also fail to compare TSCP with
any of the approaches from Multi-Modal Conformal Prediction Regions
with Simple Structures by Optimizing Convex Shape Templates
(Tumu et al.\ 2024), which can be used to form convex (including
hyper-rectangular) prediction regions.}

We appreciate this suggestion. \AppShape{} now compares TSCP with
a hyperrectangular shape-template baseline built from the public
\texttt{conformal-region-designer} components associated with
Tumu et al.\ \cite{tumu2024}. This family is relevant because it
learns geometrically simple prediction regions, including rectangular
components, rather than requiring all prediction sets to have an
unstructured boundary.

\paragraph{Configuration and implementation disclosure.}
The experiment uses two-dimensional independent Gaussian residuals
with scales $(2,1)$, $n\in\{30,50,100,200\}$, and 100 independent
repetitions. Both methods use 2,400 training and 600 test observations.
For the shape-template method, signed calibration residuals are split
equally between density/template fitting and conformal calibration.
We use KDE on a grid with 20 locations per coordinate, mean-shift
clustering, and axis-aligned hyperrectangular templates. The resulting
region may have multiple rectangular components. The comparison is
therefore with this specific multimodal template implementation, not
with a claim that every template region is one rectangle.

The implementation uses score-consistent inflation and an aspect-ratio
floor of $0.05$ for clusters that are nearly degenerate in a
coordinate. These modifications are disclosed in the appendix. We
do not present the result as an unmodified execution of every
released template class. Shape-template volume is divided by $2^d$
to match the normalized residual-volume convention used for the
other absolute-residual methods.

\paragraph{Results and scope.}
At $n=100$, the shape-template baseline has mean coverage $0.904$,
normalized volume $10.07$, and construction time $98.7$ milliseconds.
TSCP has corresponding values $0.906$, $8.04$, and $0.80$
milliseconds. The implementation's Cartesian KDE grid has $20^d$
locations, which motivates restricting this benchmark to low
dimension. This limitation is a property of the chosen computational
configuration, not a claim that every method in the shape-template
framework has exponential complexity or is unusable in higher
dimensions. Likewise, a unimodal Gaussian experiment does not test
the full potential advantage of flexible multimodal regions.

\revisionblock{R2.M5: related-work revision}{Insert the new
literature-review paragraph on Tumu et al., its distinction from
TSCP, and the reference to the low-dimensional comparison. The
published proceedings title differs slightly from the preprint title
quoted by the reviewer; use the bibliography entry below or the
corresponding entry in the current manuscript consistently.}

\reviewercomment{R2.M6. Heteroskedasticity in only some coordinates}{%
Standardizing the residuals to be used as scores is done to address
heterogeneity across dimensions. How does it work if some, but not
all, dimensions have heteroskedastic errors? A simulation
demonstrating this would make the numerical experiments more
interesting.}

We have added this experiment to the main text
(\SecAbsolute{}, \FigHeteroskedastic{}). It separates heterogeneity
across coordinates from variation of a coordinate's noise scale with
the input. The first five of ten outcomes have conditional noise
standard deviations
\begin{equation}
 \sigma_j(x_1)=(11-j)\sqrt{\frac{1+1.5x_1^2}{2.5}},
 \qquad j\leq5,
\end{equation}
where $X_1\sim\mathcal N(0,1)$; the remaining coordinates have
$\sigma_j(x_1)=11-j$. The factor $2.5$ ensures
$\E[\sigma_j^2(X_1)]=(11-j)^2$, so the comparison does not simply
increase the marginal variance of the heteroskedastic coordinates.

With $n=100$, TSCP attains mean joint coverage approximately $0.904$
and mean normalized volume $1.40\times10^{11}$. Point CHR attains
$0.902$ and $2.94\times10^{11}$, and Unscaled Max attains $0.900$
and $2.33\times10^{13}$. The coordinate profile in the same figure
identifies the heteroskedastic coordinates and shows the widths of
the final intervals. The study uses the same 200-repetition,
7,200-training/800-test protocol as the other main
absolute-residual experiments.

These observations concern joint marginal coverage over new
covariates and responses. The global residual standardization used
here does not estimate $\sigma_j(x_1)$ and does not guarantee
conditional coverage at every $x_1$. The revised experimental
discussion states this limitation explicitly. In particular,
successful average calibration should not be interpreted as having
solved covariate-dependent undercoverage in every subgroup.

\reviewercomment{R2.M7. Absolute-value CQR scores, constant shifts, and choosing the constant}{%
Above equation 9, you state one possible score is the absolute value
of CQR. It seems like this score would lead to a larger adjustment
than the standard CQR score, which cannot be used because of
Assumption 1. Is that not the case? You mentioned earlier in the
paper that a constant could be added to the standard CQR scores so
they are non-negative, and not overly conservative when the
quantile regression coverage is conservative. Would it not be
beneficial to include this score instead of the absolute value of
the CQR score, so that practitioners don't use an overly
conservative approach? Along the same line, it would be nice to
know how to choose such a constant in practice, so that this
approach could be used with quantile regression.}

We agree that the transformations should not be treated as
interchangeable. This comment motivates both the choice of score in
the main experiment and the additional shifted-score analysis.

\paragraph{Why the main experiment uses the positive part rather than
the absolute value.}
For an ordered base interval, a signed CQR score $S_j$ is negative
inside the interval. Taking $|S_j|$ assigns a large positive value
even to an observation well inside a wide base interval. Capping
instead uses
\begin{equation}
 E_j(x,\bm y)=\max\{S_j(x,\bm y),0\}.
 \label{eq:capped-score}
\end{equation}
At the score level, $E_j\leq|S_j|$ and negative signed scores are
not turned into additional evidence of prediction error. This
pointwise comparison does not by itself prove that all adaptively
standardized TSCP intervals are ordered under the two transformations.

For a nonnegative coordinate threshold $W_j$, the capped score has
the direct inversion
\begin{equation}
 E_j(x,\bm y)\leq W_j
 \quad\Longleftrightarrow\quad
 y_j\in[\widehat q_{\ell,j}(x)-W_j,
        \widehat q_{u,j}(x)+W_j].
 \label{eq:capped-inversion}
\end{equation}
Thus it produces an interval that contains the base interval. It
cannot exploit a negative adjustment to shrink an already
conservative base interval. We acknowledge this limitation explicitly,
and we leave the signed scores unchanged for the Unscaled Max,
Empirical Copula, and native CQHR comparators.

There is also a distinction in inversion: $|S_j|\leq W_j$ imposes
both $S_j\leq W_j$ and $S_j\geq-W_j$. The latter constraint can
exclude points deep inside the base interval. Consequently, the
absolute-value score does not generally have the simple expanded
interval inversion in \eqref{eq:capped-inversion}. The score and
interval formula in the revised counterpart of equation (9) must
be made consistent.

\paragraph{The constant-shift alternative is reported in the appendix.}
For a fixed $C$, the alternative uses $E_j^{(C)}=S_j+C$. If TSCP
returns threshold $W_j^{(C)}$ and nonnegativity holds on the
relevant score domain, the outcome-space adjustment is
$A_j^{(C)}=W_j^{(C)}-C$, giving
\begin{equation}
 [\widehat q_{\ell,j}(x)-A_j^{(C)},
  \widehat q_{u,j}(x)+A_j^{(C)}].
 \label{eq:shifted-inversion}
\end{equation}
An interval is interpreted as empty if its lower endpoint exceeds
its upper endpoint. Unlike capping, this inversion can in principle
allow a negative adjustment.

\AppCQR{} and \FigShift{} test $C\in\{100,300,1000\}$ with
$d=3$, $n=100$, base-interval miscoverage $0.5$, and 30 independent
repetitions. Mean joint coverage is approximately $0.9077$, and
mean outcome-space volume is approximately $2.95\times10^4$ at
each value of $C$. The recorded coordinate lengths agree exactly
between shifted TSCP and shifted TSCP-GWC within each trial at each
tested $C$. Across the three constants, the largest discrepancy in
the paired TSCP coordinate lengths is approximately
$2.2\times10^{-12}$, consistent with numerical roundoff.

These observations show that the tested shifts do not produce a
distinct local refinement in this design. They do not prove that
the practical local algorithm is translation-invariant for every
dataset or every admissible $C$. In particular, the nonnegative
domain enters the local worst-case calculation, so invariance of
an ideal standardized score under a common translation should not
be transferred automatically to all computational approximations.

\paragraph{How a shift can be justified in practice.}
For ordered quantile limits, write
$w_j(x)=\widehat q_{u,j}(x)-\widehat q_{\ell,j}(x)$. Then
\begin{equation}
 S_j(x,\bm y)\geq-\frac{w_j(x)}{2}
 \qquad\text{for every }y_j.
\end{equation}
Consequently, a known uniform bound $w_j(x)\leq M$ on the relevant
input domain gives the sufficient choice $C\geq M/2$. Such a bound
must come from the problem or the independently trained model,
not merely from the smallest observed calibration score. Choosing
$C=-\min_{i,j}S_j(X_i,\bm Y_i)$ from calibration observations
alone ensures nonnegativity only on those observations and does
not establish it for future inputs or candidate outcomes. It also
requires care about the symmetry of any data-dependent transformation.

Our fixed-constant experiment is therefore a sensitivity analysis,
not a validated general-purpose selector for $C$. We do not
recommend tuning $C$ to minimize reported test volume or to improve
test coverage. When a valid lower bound is unavailable, the capped
construction avoids that choice, while retaining the limitations
described above.

\paragraph{Ties and the revised assumptions.}
Capping produces point masses at zero. The current implementation
adds a tiny reproducible perturbation to scalar calibration
upper-bound scores, but this does not remove the zeros from the
coordinate residuals or resolve every possible zero-scale case.
It is not a proof that the original no-point-mass assumption holds.
The revised theory must state the applicable tie/degeneracy
convention and establish the claimed guarantee under it, or clearly
qualify the CQR results as empirical rather than covered verbatim
by the original assumption.

\revisionblock{R2.M7}{Insert the actual corrected CQR-score
definition, the matching inversion formula replacing equation (9),
and the revised assumptions or explanation covering ties and
degenerate scales. Insert the precise statement about a
calibration-independent shift and any verified support/width bound
used in practice. Do not assert that the experiment's calibration
nonnegativity check is a uniform guarantee over future inputs.}

\subsection{Minor Comments}

\reviewercomment{R2.m1. Marginal coverage labels in Figure 1}{%
In Figure 1 it would be a nice addition to include the marginal
coverage rates for the rightmost image. Making the marginal
coverages be different (e.g., 92\% and 93\%) would also highlight
early on that the marginal dimension coverages need not be the same.}

We agree that explicit marginal labels would make the illustration
more informative. For intervals $I_1(X)$ and $I_2(X)$, the quantities
\[
 p_1=\Prob\{Y_1\in I_1(X)\},\qquad
 p_2=\Prob\{Y_2\in I_2(X)\}
\]
need not be equal, and neither should be confused with
$\Prob\{Y_1\in I_1(X),Y_2\in I_2(X)\}$. The latter is at most
$\min(p_1,p_2)$. Thus unequal marginal coverages are compatible
with a common joint target and help convey the distinction.

\revisionblock{R2.m1: Figure 1}{Insert the updated schematic,
caption, and final figure location. Use marginal coverage values
computed from the illustrated experiment, or explicitly label
chosen values such as 92\% and 93\% as schematic. Do not attach
invented empirical percentages to an unchanged scatterplot.}

\reviewercomment{R2.m2. Meaning of ``uniformly tight'' in the population oracle}{%
In section 2.5 when discussing the oracle prediction set your
second claim is: ``it tends to be uniformly tight across all
output dimensions, because it operates with residuals whose
components are all on the same scale''. Does tight here refer
to the marginal coverage, or something else? If it is the
coverage, would the error distribution shapes need to be the
same for this to be true?}

The reviewer is right that standardization of the first two moments
does not imply identical marginal coverage. The intended point is
adaptation to coordinate scale, not a distribution-free equalization
of coordinate-wise miscoverage or a coordinate-wise optimality
claim. Standardized residual distributions may have different
skewness, tail behavior, and shapes, even when all have mean zero
and variance one.

With population parameters, a common standardized cutoff $c$
corresponds to coordinate thresholds $\mu_j+c\sigma_j$. If the
standardized residuals share a common marginal distribution, their
coverage at a fixed common cutoff is the same. However, a
finite-calibration, data-dependent cutoff requires additional care;
appropriate symmetry of the joint standardized residual law, or
an asymptotic argument under suitable regularity, is needed for
stronger balance statements. Matching means and variances alone
does not supply such a result. Cross-coordinate dependence also
affects the joint quantile.

The more precise description is that the oracle reduces
inefficiency caused by differences in location and scale and
produces coordinate-adaptive intervals. The coverage guarantee
remains joint marginal coverage; equality of marginal coverages
is not a general consequence. The new coordinate-length plots
illustrate width adaptation without claiming equal marginal
coverage across coordinates.

\revisionblock{R2.m2}{Insert the replacement for ``uniformly
tight'' in the revised Section 2.5 and any related introduction
or post-Theorem-1 claims. Define whether ``tightness'' means
interval width, volume, approximation to an oracle, or coverage,
and state extra assumptions for any balance claim that is retained.}

\reviewercomment{R2.m3. Labeling the boundaries in Figure 2}{%
In section 3.4.4, after you introduce $L_j$, you discuss what
$L_1$ and $L_2$ are with respect to Figure 2. Adding $L_1$ and
$L_2$ to the margins (or with a dashed line across the plot)
would be a helpful visualization.}

We agree. The labels should identify the horizontal extent $L_1$
and vertical extent $L_2$ of the reported residual-space rectangle,
with guide lines that distinguish these limits from the internal
partition boundaries. This directly connects the geometric
illustration to the coordinate-wise maximization defining $L_j$.
The caption should also state whether the plot is in residual
space or outcome space; for absolute residuals, $L_j$ is a
half-width and the corresponding outcome interval has full
length $2L_j$.

\revisionblock{R2.m3: Figure 2}{Insert the actual updated
figure/caption and location. Confirm that the labeled boundaries
match the displayed final rectangle and the current definition
of $L_j$.}

\reviewercomment{R2.m4. Runtime comparisons in the main text}{%
The authors discuss the computational cost of using TSCP as they
introduce it. There are some plots in the appendices which compare
the computational cost. They should either add a note in the
numerical studies that computational cost comparisons can be
found in the appendix, or (better) include similar plots in the
main text.}

We have followed the latter suggestion. The main synthetic
comparison includes a construction-time panel (\FigAbsolute{}),
and the real-data section includes a dedicated runtime plot
(\FigRuntime{}). The latter reports milliseconds on a logarithmic
axis across all six retained benchmarks. The text explicitly
states that these timings exclude model fitting and residual
computation, and that TSCP is not faster than every simpler
baseline. Additional construction-time comparisons appear in
\AppDependence{} and \AppShape{}.

As discussed in R1.C3, these measurements provide direct
wall-clock evidence for the tested sample sizes without changing
the worst-case complexity statement. The missing branch-frequency
diagnostic is addressed separately in R1.Q2 and R2.M3 rather than
being inferred from this figure.

\reviewercomment{R2.m5. Infinite Point CHR volume on the stock dataset}{%
In the stock data analysis, the volume of Point CHR is infinite
because of the size of the training/calibration sets along with
the coverage rate. One would not split their data in such a way
in practice, which should be stated. The comparison does make
the point that Point CHR does not use the data as efficiently,
but it should be clarified that the infinite prediction sets
are caused by how small the last calibration set size is, and
not the method.}

We agree, and \SecReal{} now explains this directly. The common
real-data protocol uses training/calibration/test proportions
$75\%/5\%/20\%$. On \texttt{stock}, this leaves only $n=15$
observations in the initial calibration set. Point CHR's additional
split then leaves too few observations in the final conformal
stage for a finite $0.9$ empirical conformal quantile.

More generally, with $m$ observations in the final conformal
stage, the required order-statistic index is
\[
 k=\lceil(1-\alpha)(m+1)\rceil.
\]
If $k>m$, the conformal quantile includes the augmented value
$+\infty$. At $\alpha=0.1$, a finite index requires at least
$m=9$ final-stage observations. This explains the infinite output
under the tested split; it is not an intrinsic inability of
Point CHR to produce finite regions on this dataset.

A practitioner using Point CHR would allocate more observations
to its final calibration stage, with a corresponding reduction
elsewhere in the allocation. We now state this explicitly instead
of treating the infinite-volume entry as an unconditional failure
of the baseline. The small-calibration stress test in \AppSmall{}
and \FigSmall{} further reports the frequency of infinite regions
separately from finite-volume summaries. Under that experiment's
allocation, Point CHR is infinite in every repetition at $n=10$
and finite in every repetition for the displayed $n\geq20$.

This is a comparison under a common initial data allocation, not
a comparison after optimizing the allocation separately for
every method. We have not added a separately retuned Point CHR
real-data experiment and do not claim to have done so.

\reviewercomment{R2.m6. Vector and scalar notation}{%
Readability would be improved if vectors were bolded while
scalars were not.}

We agree that a consistent typographic distinction would reduce
the burden of reading the indexed formulas. The intended
convention is to use bold symbols for vectors, such as
$\bm Y$, $\bm E$, $\bm\mu$, $\bm\sigma$, and the multi-index
$\bm h$, while leaving scalar coordinates $Y_j$, $E_j$,
$\mu_j$, $\sigma_j$, and $h_j$ unbolded. The same convention
should apply in algorithm inputs and outputs, theorem statements,
and figure captions, not only in the main equations.

\revisionblock{R2.m6}{Insert a concise confirmation of the
actual notation pass and the location of the revised notation
paragraph. Check the current manuscript and supplement
consistently; do not claim a completed global edit based only
on the experimental replacement files.}

\section{Final Consistency Checks}
\label{sec:checks}

\authoraction{This section is an author-facing completion checklist
and should be removed before submitting the response letter.}

\begin{enumerate}
 \item \textbf{Theoretical scope.} Match R1.C2, R1.Q1, and R2.M7
 to the actual revised assumptions and proofs. In particular,
 distinguish finite-sample validity from population-moment
 arguments, tie handling, nonnegative scores, and positive
 empirical scales. Do not silently upgrade a simulation to
 a theorem.
 \item \textbf{Unmeasured diagnostics.} Fill the real-data
 fallback and backward-search results in R1.Q2/R2.M3, or state
 explicitly that the request is only partially answered.
 Supply the real-data coordinate comparison in R2.M2 or
 explain its omission. None of these measurements can be
 recovered from aggregate runtime and volume alone.
 \item \textbf{CQR definitions.} The body uses
 $\max\{S_j,0\}$; the appendix reports $S_j+C$ for
 $C=100,300,1000$; native CQHR and the specified signed-score
 baselines must not be mislabeled as capped. Verify the
 base-interval inversion and disclose the reference-coordinate
 and numerical-floor choices for CQHR.
 \item \textbf{Uncertainty.} Retain the actual means, including
 those below target. The nominal target lies within one
 repetition-level standard deviation of all 58 audited
 displayed TSCP points, but one standard deviation is not
 a confidence interval for the mean. For the four reported
 synthetic shortfalls, the pointwise normal-approximation
 intervals based on $s/\sqrt B$ contain the target. This
 is a Monte Carlo diagnostic, not proof of validity or a
 simultaneous confidence statement. Real-data random-split
 variability describes resampling a fixed dataset; do not
 call those splits independent draws of new datasets.
 \item \textbf{Runtime and comparison scope.} Preserve the
 distinction between historical real-data timings and newly
 generated synthetic timings. Confirm any hardware statement
 from the run provenance. Keep the shape-template comparison
 specific to the disclosed two-dimensional implementation.
 \item \textbf{Figure and manuscript references.} Replace
 every reference macro in the preamble and each blank
 revision/location block. The existing absolute-overview
 paragraph still mentions coordinate lengths as part of
 that figure, although the dedicated bars are now in a
 separate figure; update that pointer when assembling the
 manuscript. Check that every final caption uses the
 correct full-length versus normalized-volume convention.
 \item \textbf{Presentation changes.} Insert the actual
 literature-review corrections, the modified Figures 1
 and 2, and the notation revision. Confirm that all
 modified manuscript sections are visibly highlighted
 as requested by the Editor.
 \item \textbf{Submission hygiene.} Remove author-action
 paragraphs and this checklist only after resolving the
 underlying items. Search for \texttt{INSERT},
 \texttt{pending}, \texttt{revisionblock}, and
 \texttt{authoraction}; a search hit in a macro invocation
 may still indicate unfinished content. Compile twice
 and inspect page breaks, equations, and references.
\end{enumerate}

\begin{thebibliography}{9}

\bibitem{sampson2024}
M. Sampson and K.-S. Chan (2024).
Conformal multi-target hyperrectangles.
\emph{Statistical Analysis and Data Mining: The ASA Data Science Journal},
17(5), e11710.
\url{https://doi.org/10.1002/sam.11710}.

\bibitem{tumu2024}
R. Tumu, M. Cleaveland, R. Mangharam, G. Pappas, and L. Lindemann (2024).
Multi-modal conformal prediction regions by optimizing convex shape
templates.
In \emph{Proceedings of the 6th Annual Learning for Dynamics \& Control
Conference}, Proceedings of Machine Learning Research, 242:1343--1356.
\url{https://proceedings.mlr.press/v242/tumu24a.html}.

\end{thebibliography}

\end{document}
